\FloatBarrier
\Needspace{12\baselineskip}
\section{系统设计与方法}
\label{sec:methods}
EAR 的三个工作流围绕可保存的研究状态推进：已有材料进入下一轮执行，新证据改变后续行动。数学研究在此基础上加入策略演化，让完整研究的结果参与执行做法的更新。

\FloatBarrier
\subsection{系统架构：把研究状态保留下来}
\label{sec:architecture}
一次执行从方向、资源和交付要求出发，调用工具完成检索、推导辅助、实现与评价，并留下研究材料和反馈。下一轮据此继续工作。图\ref{fig:overview} 展示提案、稿件与评阅材料之间的关系；三个模块可独立接收输入，跨阶段的材料交接由研究者组织。本文的效果证据来自各模块的独立记录，跨模块完整流程尚未单独评估。

\DiagramFig{\resizebox{.97\textwidth}{!}{\SystemOverview}}{EAR 的研究任务与产物衔接。三个工作流分别可用，提案、稿件和评阅材料可以成为后续输入；关键判断与研究责任由人承担。}{fig:overview}

\begin{table}[H]
\centering\small
\begin{tabularx}{\textwidth}{P{24mm}Y Y}
\toprule
\rowcolor{warm}
\textbf{工作流} & \textbf{输入与交付} & \textbf{人保留的判断}\\
\midrule
问题发现 & 方向与条件 $\to$ 文献比较、候选、提案 & 问题价值与执行路径\\
数学研究 & 方向或提案 $\to$ 论证、代码、稿件与评价 & 关键技术结论与验收\\
评审回应 & 论文、评审与证据 $\to$ 回应及证据账本 & 主张边界与最终责任\\
\bottomrule
\end{tabularx}
\caption{三个工作流的输入、交付与人工职责。研究者通过提案、稿件和评阅材料组织阶段交接。}
\label{tab:interfaces}
\end{table}

\Sub{交接做过的判断，研究才能继续}
研究中途换一个执行者，最需要交接的是做过的判断：采用了哪些设定，哪些文献最接近，证明推进到了哪里，哪条路线已经失败。EAR 保存这些状态及后续验证计划，使工作能够从已有结果继续。评价记录同时标明实际工具反馈与宿主自审的来源，研究者可以据此检查重要取舍。\src{1}

\FloatBarrier
\subsection{问题发现：批判性树搜索，让问题逐轮具体化}
\label{sec:idea}
问题发现从研究方向与资源条件起步，把候选逐步展开为明确的问题、假设、方法和验证计划。每次扩展都引入新的检索与批判，节点的价值随证据更新，下一次投入也随之调整。图\ref{fig:tree} 给出候选扩展和分支退出的关系。

\DiagramFig{\resizebox{.96\textwidth}{!}{\CriticalTree}}{候选扩展与批判性搜索。检索和评价更新节点价值，碰撞、不可行或重复的分支可以退出，退出依据仍然保存。}{fig:tree}

\Needspace{7\baselineskip}
\subsubsection{每次扩展，都重新检索与批判}
当前实现用 UCT 引导 best-first 扩展。模型提出子节点，评价工具更新访问次数与节点价值，并沿父链回传。每个新候选至少重新检查新颖性和研究适配度；选择节点时，综合评分、可核查的文献重合、验证规则与结构信号共同参与判断。搜索直接评价展开后的候选，不使用随机 rollout 模拟终局。\src{1,2}

表\ref{tab:idea-scores} 列出四类评价及组合权重。分数与理由、检索来源一起保存，使继续扩展某个分支的依据可追溯。不同 Agent 承担生成、批判和评价职责；评阅者画像的调用范围，以实际提示和执行记录为准。

\begin{table}[H]
\centering\small
\begin{tabularx}{\textwidth}{P{25mm}Y P{14mm}}
\toprule
\rowcolor{warm}
\textbf{字段} & \textbf{检查的研究问题} & \textbf{权重}\\
\midrule
Novelty & 最近工作与核心机制是否具有实质差异 & .25\\
Venue & 研究价值、严谨性与目标要求是否匹配 & .35\\
Strategic & 是否可证伪、证据可得，是否符合资源条件 & .20\\
Feasibility & 实现复杂度、算力、数据与时间线是否可行 & .20\\
\bottomrule
\end{tabularx}
\caption{当前搜索评价字段与权重。评价结果影响下一次扩展，来源和理由为复核提供依据。}
\label{tab:idea-scores}
\end{table}

\Needspace{7\baselineskip}
\subsubsection{从研究意图到可执行提案}
研究者给出方向、资源约束与关键追问，Claude Code Harness 组织任务并调度工具，Codex MCP 返回候选、批判、评分与建议。图\ref{fig:harness} 展示选题模块的一种交互接入方式：反馈返回当前任务，筛选与精修结果保存下来，成为下一步工作的输入。

\DiagramFig{\resizebox{.96\textwidth}{!}{\HarnessAccess}}{选题模块的交互式接入。研究者给出约束并选择下一步，Harness 与评价工具交换反馈，保存候选、文献和精修结果。}{fig:harness}

候选筛选与提案精修可以分别进行。筛选时，评价工具读取候选的核心主张与近邻文献，判断新颖性和实质差异，并给出继续或放弃的建议。精修阶段沿保留的候选展开，补齐研究设定与验证计划。早期接入以分阶段交互推进，当前树搜索进一步组织了候选扩展；两者保留了同样重要的材料交接。\src{1}

\Needspace{7\baselineskip}
\subsubsection{退出理由也是研究产物}
候选与已有工作实质重合，或在给定条件下不可行时，搜索可以停止该分支。EAR 保留节点、来源、差异分析与退出理由，后续工作仍能利用这部分负向证据。留下的方案则围绕确定的问题锚点继续精修，明确假设、主要论据、实现接口和最低验证计划，形成可执行提案。

\begin{insight}
\textbf{过程案例 E1：新增近似工作，使高分候选退出}\par
在一次早期分阶段复判中，补充检索在两个候选周围共找到三项近似工作。候选 A 的新颖性由 $8/10$ 降至 $4/10$，继续研究的建议转为放弃；候选 B 由 $8/10$ 降至 $5/10$。转折来自新增文献：系统重新比较核心差异，随之修改分数、理由和投入建议。\src{10}
\end{insight}

最初的检索未充分覆盖近邻工作，候选因而显得更有差异。补充检索改变了这个判断，也改变了后续行动。这段记录展示了复筛如何影响投入决定。它来自早期分阶段流程，与当前 UCT 搜索版本不同，算法间的性能比较仍需相应对照。

\FloatBarrier
\subsection{数学研究：内循环做研究，外循环改进做法}
\label{sec:math}
数学与算法研究可以把技术疑虑落实到具体命题、反例、构造和数值检查中。EAR 用内循环沿这些线索推进当前任务，再把一次完整研究的结果交给外循环，检验采用的执行策略。图\ref{fig:dual-loop} 展示研究修订与策略演化之间的嵌套关系。

\DiagramFig{\resizebox{.96\textwidth}{!}{\MathDualLoop}}{数学研究的双循环。内层推进当前任务，外层依据完整研究结果比较与更新策略，并保留初始策略。历史策略可自行寻找具体目标。}{fig:dual-loop}

\Needspace{7\baselineskip}
\subsubsection{保留研究状态，沿具体缺口修订}
内循环围绕明确目标开展选题、推理、证明或构造、实现、检查与写作。精确命题、失败的证明步骤、反例、数值验证与实验对照都进入研究状态。一次尝试结束后，下一步可以沿尚未解决的缺口继续；达到交付条件时完成稿件，无法继续时记录退出。

评阅指出证明缺口，修订就返回推导；实现或验证存在问题，则补充技术检查；贡献不足需要重新审查范围与主张，表达问题进入稿件重写。历史运行在初评后执行三轮修订与再评阅，并恢复平均评分最高的稿件版本。每轮修改及评价变化都被保留，当前轮轨迹与最终交付版本因此分别统计，评价协议见第\ref{sec:protocol} 节。

\Needspace{7\baselineskip}
\subsubsection{策略要经受一次完整研究}
外循环的基因组表示一套候选研究策略。历史版本由阶段提示与策略参数构成，涵盖目标筛选、证明组织、写作和修订。当前可移植版本将可修改范围扩展到任务程序、提示、技能、工具流程与配置，基础模型权重保持固定。初始策略与候选策略分别保存，谱系、研究结果和失败原因构成后续选择的依据。\src{1,4}

一个策略是否值得采用，要由它执行研究的结果回答。每个候选都要尝试完整研究，目标选择、技术进展、稿件与评阅结果，以及失败出口，都进入外循环选择。在当前可移植版本中，任务列表、终审与标准固定在候选的修改范围之外；历史实验允许策略自行选择具体研究目标。第\ref{sec:protocol} 节按这一历史设定解释结果。

策略更新采用四种操作。\textbf{精英保留}延续表现较好的策略；\textbf{定向变异}针对退出原因、技术缺口或低分意见修改父策略；\textbf{语义交叉}组合两个父策略的互补指令；\textbf{探索}尝试较大变化。历史改动具体落在选题与论证上：缩短候选清单，提高目标与贡献要求，并提前要求核心引理和可行证明路线。更新后的策略再次执行研究，接受结果检验。

\Needspace{7\baselineskip}
\subsubsection{成功与退出，都进入适应度}
历史 episode 的适应度按下式计算：
\begin{equation}
F(e)=\begin{cases}
q_{\mathrm{best}}(e),&\text{完成稿件且具有有效评分},\\
0.5,&\text{研究退出},\\
0,&\text{筛选淘汰},\\
-0.5,&\text{没有目标或搜索失败},\\
-2,&\text{预算耗尽或贡献过于平凡}.
\end{cases}
\label{eq:historical-fitness}
\end{equation}
$q_{\mathrm{best}}$ 为初评与修订中保留的最高稿件均分。稿件完成但所有评分尝试都失败时，历史实现采用 3.0 回退值。研究退出、筛选淘汰、搜索失败与预算耗尽保留不同适应度，使成功写作之前的取舍与失败也进入策略评价。

当轮策略按平均 $F$ 排序选择。历史精英档案另用以下组合指标保存策略表现：
\begin{equation}
A_{\mathrm{hist}}=\frac{1}{3}\left(\frac{\overline F+2}{12}+w+a\right),
\label{eq:archive-fitness}
\end{equation}
其中 $w$ 为稿件完成率，$a$ 为历史模型评阅通过率。当前可移植版本等权组合归一化质量、完成率与固定终审通过率：
\begin{equation}
A_{\mathrm{portable}}=\frac{1}{3}\left(\frac{\overline q_{\mathrm{clip}}}{10}+w+a_{\mathrm{final}}\right).
\label{eq:portable-fitness}
\end{equation}
当前质量分数按约定规则裁剪并归一化；组合得分相当时，再比较调用与返工。平均 $F$ 用于历史当轮选择，$A_{\mathrm{hist}}$ 用于历史档案，$A_{\mathrm{portable}}$ 用于当前版本。调用效率另按执行记录核算。\src{1,10}

\FloatBarrier
\subsection{评审回应：先处理关键疑虑，再组织回应}
\label{sec:rebuttal}
评审回应从论文、审稿意见、已有证据与截止时间出发。它先判断疑虑对应的证明或实验是否已经存在，再安排必要的补充工作，组织回应草稿并保留依据。图\ref{fig:rebuttal-flow} 展示从疑虑到证据、再到回答的过程。

\DiagramFig{\resizebox{.96\textwidth}{!}{\RebuttalFlow}}{从论文和评审定位疑虑及证据。已有证据足够时直接起草；有技术缺口时先补齐。模型反馈用于修订，作者审阅关键内容与最终回应。}{fig:rebuttal-flow}

\Needspace{7\baselineskip}
\subsubsection{先找到疑虑对应的证据}
几位审稿人可能从不同角度提出同一疑虑。EAR 先归并相关意见，找到论文中的对应论据，再分辨问题出在解释、证明还是实验。解释不足时，把已有证据组织清楚；确实存在技术缺口时，再安排补充研究。同一项证明或实验可以回答多条相关意见，证据账本记录主张与出处。\src{1,5,9}

社区公开的回应方法与作者经验提供初始策略。候选在回答顺序、证据选择与组织方式上有所不同，优先处理影响论文判断的疑虑。资源约束下的目标写为：
\Needspace{8\baselineskip}
\begin{equation}
\max_{\theta}\ \mathbb E[\Delta R\mid\mathrm{Context},\theta]
\quad\text{subject to}\quad H\le H_B,\ C\le B,\ T\le D.
\label{eq:rebuttal-objective}
\end{equation}
$\mathrm{Context}$ 包含论文、评审与已有证据，$\theta$ 为回应策略，$\Delta R$ 为希望改善的评分，$H_B$ 为人时上限，$D$ 为截止时间。该式表达设计目标；现有预测测试检查反馈组件，过程记录检查策略调整，生成回应对真实评分的影响尚未直接测量。

\Needspace{7\baselineskip}
\subsubsection{保留好草稿，追踪未解疑虑}
候选策略可以优先说明已有证据，围绕关键疑虑组织论证，或先补齐缺失验证。它们使用同一份论文与审稿材料，把工作安排与回答方式作为比较对象。图\ref{fig:rebuttal-population} 展示草稿和反馈怎样进入下一轮。

\DiagramFig{\resizebox{.96\textwidth}{!}{\RebuttalPopulation}}{回应策略从社区方法与作者经验初始化。候选共享材料，模型评价与事实检查形成反馈；下一轮继承较好草稿、证据、未回答的问题和修改建议。}{fig:rebuttal-population}

较好的草稿保留下来，未解疑虑、证据缺口与修改建议随它进入下一轮。这是一次黑盒优化：系统比较候选回应，读取诊断反馈后调整策略，无需计算评分对策略的导数。评价同时检查疑虑覆盖和论据依据，并预测评分上涨、不变或下降；预测能力在第\ref{sec:rebuttal-results} 节单独测试。

\begin{insight}
\textbf{过程案例：两轮反馈带来策略调整}\par
一份匿名历史记录每轮比较三种策略：优先依据已有证据、围绕关键疑虑组织论证，以及分阶段补证据。第一轮所选的证据优先策略未达到历史内部代理的停止要求；第二轮继承草稿与修改建议，改用关键疑虑优先策略后触发停止条件。记录展示了反馈如何改变写法，尚未用当前门槛重新评阅，也未测量真人评分变化。\src{10}
\end{insight}

历史流程在草稿通过内容评价后停止本轮其余策略尝试，再核对回应是否忠实于论文与证据、表述是否准确。检查发现的问题返回修订，最终草稿交给作者审阅。新增技术结论、证据有效性与最终回应由作者确认。
